\section{A bounded ultrapower test of the remaining cote 134-8 tensor gate}
\label{sec:arch134-ultrapower}

The remaining original question is the \emph{Thm(?)} on cote~134-8,
physical p.~49, footer~48, scan~049, with the tensor and interval
predicates on physical p.~48, footer~47, scan~048.  ARCH-134-02 has
already refuted automatic higher surjectivity of F-SM-44, under all
those predicates.  It has not refuted this localization theorem.
This section tests a different exact functor in the inherited
F-SM-34 additive construction, retaining its original scalar
embedding.  It does not assert that every possible tensor on that
category has been excluded.

\subsection{ARCH-134-05: the ultrapower quotient retains the precise tensor obstructions}

\paragraph{Statement.}
Let $k$ be a finite field and let $\mathcal U$ be a nonprincipal
ultrafilter on $\mathbb N$.  Define
\[
 F(V)=V^{\mathbb N}/\mathcal U,\qquad
 \eta_V:V\longrightarrow F(V),\qquad
 E(V)=F(V)/\eta_V(V),
 \tag{ARCH134.U1}
\]
where two sequences are identified when they agree on a set in
$\mathcal U$, and $\eta$ is the constant-sequence map.  Then $F$
and $E$ are exact $k$-linear functors, $F(k)=k$, $E$ vanishes on
finite-dimensional spaces, and $E(k^{(\mathbb N)})\ne0$.
The pointwise maps
\[
 \mu_{V,W}:F(V)\otimes_kF(W)\longrightarrow F(V\otimes_kW)
 \tag{ARCH134.U2}
\]
give $F$ a symmetric lax monoidal structure with its stated scalar
unit.  The following four limitations hold.

\begin{enumerate}
\item The maps (ARCH134.U2) are not generally onto.
\item They do not descend through the two quotient maps to maps
      $E(V)\otimes E(W)\to E(V\otimes W)$.
\item There is no natural linear retraction $F\to\mathrm{id}$
      whose value at $k$ is the identity.
\item There is no natural family
      $c_{V,W}:E(V\otimes W)\to E(V)\otimes W$ whose value at
      $W=k$ is the identity.
\end{enumerate}

\paragraph{Proof.}
All operations on $F(V)$ are performed on representative
sequences; the ultrafilter equivalence respects every finite linear
operation.  Every exact sequence of vector spaces splits, and
applying the additive functor $F$ to a splitting proves exactness.
The map $\eta_V$ is injective: a nonzero constant sequence is not
zero on any member of $\mathcal U$.  Apply the snake lemma to the
map from an exact vector-space sequence to its exact $F$-image.
The three injective vertical maps show that the three cokernels
form an exact sequence.  Thus $E$ is exact as well.

Every map from $\mathbb N$ to a finite set has exactly one fiber
in $\mathcal U$.  Hence every sequence in a finite-dimensional
$k$-space is equal, modulo $\mathcal U$, to a unique constant
sequence.  This proves $F(V)=V$ and $E(V)=0$ in finite dimension.
In contrast, put $V=\bigoplus_{n\ge0}k v_n$.  The sequence
$(v_n)$ is not equal to any fixed vector on a member of
$\mathcal U$: each such equality holds at most once.  Consequently
\[
                  e=[(v_n)]\bmod\eta_V(V)\ne0.
 \tag{ARCH134.U3}
\]

The formula for (ARCH134.U2) is
$[(x_n)]\otimes[(y_n)]\mapsto[(x_n\otimes y_n)]$.
It is well defined and bilinear because intersections of finitely
many members of $\mathcal U$ remain in $\mathcal U$.
Associativity, symmetry, and both unit identities are the
corresponding identities at every coordinate.  The unit is
$k\xrightarrow{\eta_k}F(k)=k$.

For failure of surjectivity, take another countable space
$W=\bigoplus_{n\ge0}kw_n$ and consider
\[
           z_n=\sum_{i=0}^n v_i\otimes w_i.
 \tag{ARCH134.U4}
\]
Each $z_n$ has tensor rank $n+1$.  An element in the image of
(ARCH134.U2) is a finite sum of, say, $r$ pure tensors of sequence
classes, so it has a representative of rank at most $r$ at every
coordinate.  Equality to $[(z_n)]$ would force $n+1\le r$ on a
member of $\mathcal U$, impossible since that set is finite.

For the quotient failure, set $V=k$ and retain a class
$\xi\in F(W)$ with nonzero quotient $\bar\xi\in E(W)$.
The pair $(\eta_k(1),\xi)$ has zero image in $E(k)\otimes E(W)$,
whereas its product has quotient $\bar\xi$ in
$E(k\otimes W)=E(W)$.  Thus even this one pair prevents descent.
Equivalently, a tensor ideal in the pointwise structure which
contains the whole constant scalar unit contains every element;
the quotient in (ARCH134.U1) is not such a monoidal quotient.

For the retraction claim, suppose $\rho:F\to\mathrm{id}$ is
natural and $\rho_k=1$.  Use $V$ and $\xi=[(v_n)]$ above.
For each coordinate functional $\epsilon_j:V\to k$,
$F(\epsilon_j)(\xi)=0$, since its representative is supported
at the single index $j$.  Naturality forces every coordinate of
$\rho_V(\xi)$ to be zero; hence $\rho_V(\xi)=0$.  The linear
functional $\lambda:V\to k$ defined by $\lambda(v_n)=1$ for
every $n$ instead satisfies $F(\lambda)(\xi)=1$.
Naturality then gives
$\lambda\rho_V(\xi)=\rho_k(1)=1$, a contradiction.

Finally use the same $V,W$ and define
\[
 u=\left[(v_n\otimes(w_0+\cdots+w_n))_{n\ge0}\right]
           \bmod\eta_{V\otimes W}(V\otimes W).
 \tag{ARCH134.U5}
\]
For the $j$-th coordinate functional $\epsilon_j$ on $W$, the
sequence obtained by applying $1\otimes\epsilon_j$ equals
$(v_n)$ for every $n\ge j$.  Consequently
\[
                 E(1\otimes\epsilon_j)(u)=e
                       \quad\hbox{for every }j.
 \tag{ARCH134.U6}
\]
Naturality of $c$ and its unit normalization would give
$(1\otimes\epsilon_j)c_{V,W}(u)=e$ for all $j$.  Every element
of $E(V)\otimes k^{(\mathbb N)}$ has only finitely many nonzero
coordinates, contradicting (ARCH134.U3).  This proves all claims.
\hfill$\square$

\paragraph{Exact consequence for F-SM-34.}
The inherited two-sort category and projective resolution use
only exact $k$-linearity of $E$:
\[
 (X_0,X_1,b:X_0\to X_1,a:E(X_1)\to X_0),\qquad
 ba=0,\quad aE(b)=0,\quad \kappa(V)=(V,0).
\]
Their formulas therefore still give
$L_{2r}\varphi\kappa(V)=E^r(V)$ and
$L_{2r+1}\varphi\kappa(V)=0$ for the functor (ARCH134.U1).
In particular, the additive scalar defect remains and the global
unit has no positive derived defect.  This is an application of
the supplied F-SM-34 calculation, not a second proof of it.

The complete inherited proofs F-SM-36--38 and F-SM-41 were read
again for this application.  The new element (ARCH134.U5)
replaces the double-dual element in F-SM-36.  Hence its componentwise
scalar-tensor obstruction, and the F-SM-38 two-extension obstruction
to exact tensoring with $\kappa(k^{(\mathbb N)})$, persist unchanged.
F-SM-37 already treats every nonzero exact additive $E$, so its
arrow-monad strength contradiction applies directly.  The
composition contradiction of F-SM-41 uses only a nonzero
$e\in E(W)$ and the normalized forward strength, and also applies.
None of these conclusions prohibits a different, nonexact tensor
with enlarged underlying components.  Nor does the forward map
$E(V)\otimes W\to E(V\otimes W)$ provide the missing reverse
map required by the two-sort operation $a$.

\paragraph{Primary-source scope.}
Coulembier, \emph{Homological kernels of monoidal functors},
Selecta Mathematica~29 (2023), article~24,
\S6.2.3, discusses a symmetric monoidal ultrapower category and
the identification of the scalar ultrapower with $k$ for finite
$k$.  That categorical construction does not identify the
quotient $F/\mathrm{id}$ above with a monoidal quotient.
The elementary arguments (ARCH134.U1)--(ARCH134.U6) establish
the precise claims used here; no assertion about an abelian
envelope is imported.
\url{https://doi.org/10.1007/s00029-023-00829-y}

\paragraph{Verification status.}
The author completed a whole saved-proof reading: exactness and
finite-field normalization, the unbounded tensor-rank witness, the
scalar-unit cross term, both naturality contradictions, and the
exact inherited tensor-obstruction scopes were checked.  No gap
was found. The root also completed a whole saved-proof reading;
the actual events are recorded in the register.  The
exact conclusion is a failed specified tensor route, not a
counterexample to the full sheet.
